\documentclass[10pt,a4paper]{amsart}
\usepackage[margin=19mm]{geometry}
\usepackage{amsmath,amssymb,mathtools}
\usepackage[scr=rsfs,cal=euler]{mathalfa}
\usepackage{xfrac}
\usepackage[unicode,hidelinks,bookmarks=false]{hyperref}
\newcommand{\ie}{i.e.\ }
\newcommand{\cf}{cf.\ }
\newcommand{\resp}{respectively\ }
\newcommand{\Subsection}{\subsection}
\providecommand{\nobreakdash}{\nobreak\hbox{-}}
\setlength{\emergencystretch}{2em}
\multlinegap0pt
\usepackage{bm}
\usepackage{bbm}
\usepackage{dsfont}
\usepackage{xspace}
\usepackage{cancel}
\usepackage{mathtools}
\usepackage{amsthm}
\makeatletter
\@ifundefined{theorem}{\newtheorem{theorem}{Theorem}}{}
\@ifundefined{lemma}{\newtheorem{lemma}[theorem]{Lemma}}{}
\makeatother
\title[Concentration, Local Uniqueness, and Morse Index of Multi-bu]{Concentration, Local Uniqueness, and Morse Index of Multi-bubble Solutions for a Critical Exponential Biharmonic Choquard Equation}
\author{Wenjing Chen, Shengbing Deng}
\date{Source: arXiv:2608.03321v1 (CC BY 4.0)}
\begin{document}
\maketitle

\noindent\emph{Source and licence.} Concentration, Local Uniqueness, and Morse Index of Multi-bubble Solutions for a Critical Exponential Biharmonic Choquard Equation. Version arXiv:2608.03321v1 (CC BY 4.0), distributed under the Creative Commons Attribution 4.0 International licence, as recorded in the arXiv licence metadata for this exact version and shown on the abstract page https://arxiv.org/abs/2608.03321v1. Full rights notice: licenses/arxiv-2608.03321-CC-BY-4.0.txt.\par\smallskip

\noindent\textit{Editorial scope.} Counted unit: the Lemma whose source label is \texttt{lem:quadratic-localization} in the article (source0.tex lines 3685--3733, source label \texttt{lem:quadratic-localization}), delivered together with the article's own complete original proof of it (source0.tex lines 3735--3961, one \texttt{proof} environment of 227 lines). No mathematics is written, repaired or re-derived: statement and proof are the article's own text line for line. Required context retained uncounted: lem:weighted-low-mode-control (lines 763--778); lem:separated-core-estimates (lines 4310--4329). Every definition, display and label of those lines is retained unchanged. Omitted from this package and not redistributed: 1--762, 779--3684, 3734--3734, 3962--4309, 4330--4849. Nothing inside a retained excerpt is omitted or altered: every retained body is delivered exactly as the article's lines are written, so no omission receipt of any kind accompanies this package. Citations: the retained excerpts carry no citation command at all, so no bibliography entry of the article is transcribed and no reference is redistributed. Reference adaptation: none was needed and none was made; every reference inside the delivered unit resolves inside it, so no unresolved reference or citation is produced. Printed slips kept literal: no printed word, symbol, hypothesis or number of the counted statement or of its proof was altered. Layout and class: the article's own class is \texttt{article} with packages including amsmath, amssymb, amsthm, mathrsfs, cite, geometry; the wrapper is the lane's pinned \texttt{amsart} editorial wrapper at 10pt on A4, which restates the theorem-like environments the retained text needs and adds the article's own macro definitions that the retained text needs (none), with extra wrapper packages (bm, bbm, dsfont, xspace, cancel, mathtools). The delivered numbering is the unit's own. Assets: no retained span contains a figure environment, a graphics-inclusion command, a PDF-page inclusion, a table or a TikZ picture; no image file is copied into this package, the package declares no asset dependency and no image file is redistributed with it. Licence: version arXiv:2608.03321v1 of this article is distributed under the Creative Commons Attribution 4.0 International licence, as recorded in the arXiv licence metadata for this exact version and shown on the abstract page https://arxiv.org/abs/2608.03321v1. A full rights notice is delivered at licenses/arxiv-2608.03321-CC-BY-4.0.txt. Characters: every retained excerpt is pure ASCII, so the delivered engine, XeLaTeX with its default Latin Modern text font, reports no missing character.
\begin{lemma}\label{lem:weighted-low-mode-control}
There is a constant $C>0$ such that every
$\phi\in\mathcal E_\alpha$ satisfies
\begin{equation}\label{eq:weighted-low-mode-control}
\begin{aligned}
\int_{\mathbb R^4}F\phi^2\,dx
\leq C\Bigg[
&\int_{\mathbb R^4}|\Delta\phi|^2\,dx
+\left|\int_{\mathbb R^4}F\phi\,dx\right|^2+\sum_{k=0}^{4}
\left|\int_{\mathbb R^4}FZ_k\phi\,dx\right|^2
\Bigg].
\end{aligned}
\end{equation}
The same estimate holds for every translated and dilated bubble, with
the same constant after the natural change of variables.
\end{lemma}

\begin{lemma}\label{lem:separated-core-estimates}
If $i\ne j$ and $x\in B_{\eta/4}(\xi_i)$, then
\begin{equation*}
\int_{B_{\eta/4}(\xi_j)}
\frac{K(y)e^{W(y)}}{|x-y|^\alpha}\,dy
\leq C\mu_j^{4-\alpha}.
\end{equation*}
The self contribution in the $i$-th core is comparable with
\[
\mu_i^4(1+\mu_i|x-\xi_i|)^{-\alpha}.
\]
Consequently the relative cross error is $O(\varepsilon^\alpha)$ on bounded core scales. Moreover,
\begin{equation}\label{eq:separated-double}
\varepsilon^{8-\alpha}
\int_{B_{\eta/4}(\xi_i)}
\int_{B_{\eta/4}(\xi_j)}
\frac{K(x)K(y)e^{W(x)+W(y)}}{|x-y|^\alpha}\,dxdy
=O(\varepsilon^\alpha).
\end{equation}
\end{lemma}

\begin{lemma}\label{lem:quadratic-localization}
Let $\varepsilon_n\to0$ and let
$v_n\in\mathcal Y_{\varepsilon_n}$ satisfy
$\|v_n\|_{\mathcal H}\leq1$. For every large $R$, there are disjoint
logarithmic cutoffs $\eta_{i,n,R}$ such that
\[
\eta_{i,n,R}=1
\quad\text{in }B_{R/\mu_{i,n}}(\xi_{i,n}),
\qquad
\eta_{i,n,R}=0
\quad\text{outside }B_{R^2/\mu_{i,n}}(\xi_{i,n}),
\]
and
\begin{equation}\label{eq:quadratic-localization-cutoff}
|\nabla^k\eta_{i,n,R}(x)|
\leq
\frac{C}{|\log R|\,|x-\xi_{i,n}|^k},
\qquad k=1,2.
\end{equation}
Put
\[
v_{i,n}=\eta_{i,n,R}v_n,
\qquad
v_{0,n}=v_n-\sum_{i=1}^m v_{i,n}.
\]
Then
\begin{equation}\label{eq:quadratic-localization}
\begin{aligned}
\mathcal Q_{\varepsilon_n}(v_n,v_n)
\geq{}&
\sum_{i=1}^m
\mathcal Q_{\varepsilon_n}(v_{i,n},v_{i,n})
+(1-\delta_R)\|v_{0,n}\|_{\mathcal H}^2 -\delta_R\|v_n\|_{\mathcal H}^2-o_n(1),
\end{aligned}
\end{equation}
where $\delta_R\to0$ as $R\to\infty$. Moreover, after rescaling the
$i$-th core,
\[
\widetilde v_{i,n}(z)
:=v_{i,n}\left(\xi_{i,n}+\frac z{\mu_{i,n}}\right),
\]
one has
\begin{equation}\label{eq:quadratic-core-limit}
\mathcal Q_{\varepsilon_n}(v_{i,n},v_{i,n})
=
\mathcal Q_\infty(\widetilde v_{i,n},\widetilde v_{i,n})
+o_n(1)+o_R(1).
\end{equation}
\end{lemma}

\begin{proof}
We first localize the biharmonic energy. In the cylinder associated
with the $i$-th center, set
\[
s=\log(\mu_{i,n}|x-\xi_{i,n}|),
\qquad
\theta=\frac{x-\xi_{i,n}}{|x-\xi_{i,n}|},
\qquad
\widetilde v(s,\theta)=v(x).
\]
On every annulus,
\begin{equation}\label{eq:Emden-Fowler-biharmonic}
\int |\Delta v|^2\,dx
=\int
\left|\left(\partial_s^2+2\partial_s
+\Delta_{\mathbb S^3}\right)\widetilde v\right|^2
\,dsd\theta.
\end{equation}
Choose a fixed smooth function $\zeta$ with $\zeta=1$ on
$(-\infty,0]$ and $\zeta=0$ on $[1,\infty)$, and define
\[
\eta_{i,n,R}(x)
=\zeta\left(
\frac{\log(\mu_{i,n}|x-\xi_{i,n}|)-\log R}
{\log R}
\right).
\]
This gives \eqref{eq:quadratic-localization-cutoff}. For large $n$ the
supports are disjoint.

We record the one-dimensional estimate used for the logarithmic
cutoff. Expand in an orthonormal spherical-harmonic basis. We suppress
the multiplicity index and write
$\widetilde v_n=\sum_{k\geq0}a_{k,n}(s)Y_k(\theta)$. Put
$P_k=\partial_s^2+2\partial_s-k(k+2)$. Direct expansion gives
\[
P_k(\eta a)=\eta P_ka+2\eta'a'
+(\eta''+2\eta')a.
\]
For $k\geq2$, the roots of the characteristic polynomial are
separated from the imaginary axis uniformly in $k$. Integration by
parts therefore gives
\begin{equation}\label{eq:high-mode-logarithmic-cutoff}
\begin{aligned}
&\|P_k(\eta a_{k,n})\|_{L^2}^2
+\|P_k((1-\eta)a_{k,n})\|_{L^2}^2\leq
\|P_ka_{k,n}\|_{L^2}^2
+\frac{C}{\log R}\|P_ka_{k,n}\|_{L^2}^2.
\end{aligned}
\end{equation}

We treat $k=0,1$ separately.  The homogeneous solutions of the
pure biharmonic operators are
\[
1,\ e^{-2s}\quad(k=0),
\qquad
e^s,\ e^{-3s}\quad(k=1).
\]
They must not be confused with the conformal kernel of
$\mathcal L_\infty$.  Put
\[
I_R=(\log R,2\log R),
\qquad
I_R^*=\left(\frac12\log R,\frac52\log R\right).
\]
The following one-dimensional estimate is the low-mode substitute
for the spectral gap.  For $k=1$, the four multiplicity indices are
included in all sums:
\begin{equation}\label{eq:low-mode-variation-bound}
\begin{aligned}
&\sum_{k=0}^1\int_{I_R}
\bigl(|a_{k,n}'|^2+(1+k)^2|a_{k,n}|^2\bigr)\,ds\\
&\leq C\log R\Bigg(
\sum_{k=0}^1\|P_ka_{k,n}\|_{L^2(I_R^*)}^2
+|\mathfrak m_{*,n}|^2
+\sum_{\nu=0}^4|\mathfrak m_{\nu,n}|^2
+o_R(1)\|v_n\|_{\mathcal H}^2\Bigg)+o_n(1).
\end{aligned}
\end{equation}
where
\[
\mathfrak m_{*,n}=\int_\Omega
\chi_iF_{\mu_{i,n},\xi_{i,n}}v_n,
\qquad
\mathfrak m_{\nu,n}=\int_\Omega
\chi_iF_{\mu_{i,n},\xi_{i,n}}Z_{i,\nu}v_n,
\quad0\leq\nu\leq4.
\]

For completeness, we justify this estimate.  Variation of constants
on $I_R^*$ writes
\[
\begin{aligned}
a_{0,n}(s)
&=A_{0,n}+B_{0,n}e^{-2s}
+\int_{I_R^*}K_0(s,\sigma)P_0a_{0,n}(\sigma)\,d\sigma+r_{0,n}(s),\\
a_{1,n,\ell}(s)
&=A_{1,n,\ell}e^s+B_{1,n,\ell}e^{-3s}
+\int_{I_R^*}K_1(s,\sigma)P_1a_{1,n,\ell}(\sigma)\,d\sigma
+r_{1,n,\ell}(s),
\qquad 1\leq\ell\leq4,
\end{aligned}
\]
where the remainders are determined by the energy on the two
adjacent cylinder strips.  The Green kernels satisfy
\[
|K_0|+|\partial_sK_0|\leq C(1+|s-\sigma|),
\qquad
|K_1|+|\partial_sK_1|
\leq Ce^{-c|s-\sigma|}
\]
after the growing $e^s$ component has been separated.  The
coefficients of $e^{-2s}$ and $e^{-3s}$ are controlled by the
$H^2$-energy on the inner strip.  The coefficient of $e^s$ is
$O(\mu_{i,n}^{-1})\|v_n\|_{\mathcal H}$, because in the original
variables it is an affine harmonic function and
$\mathcal H\hookrightarrow H_0^1(\Omega)$.  Finally, evaluating the
constant and the remaining degree-zero and degree-one components
against the six weighted test functions gives
\[
|A_{0,n}|+\sum_{\ell=1}^4|A_{1,n,\ell}|
\leq C\left(
|\mathfrak m_{*,n}|+\sum_{\nu=0}^4|\mathfrak m_{\nu,n}|
+\sum_{k=0}^1\|P_ka_{k,n}\|_{L^2(I_R^*)}\right)
+o_R(1)\|v_n\|_{\mathcal H}+o_n(1).
\]
Here Lemma \ref{lem:weighted-low-mode-control} is used to pass from
the weighted core norms to the coefficients.  Substitution in the
two variation-of-constants formulas and integration over $I_R$
proves \eqref{eq:low-mode-variation-bound}.  The factor $\log R$ is
unavoidable for the constant plateau.

For $v_n\in\mathcal Y_{\varepsilon_n}$ the mass moment
$\mathfrak m_{*,n}$ and all the moments $\mathfrak m_{\nu,n}$ vanish. Since
$|\eta'|\leq C/(\log R)$ and
$|\eta''|\leq C/(\log R)^2$, equations
\eqref{eq:high-mode-logarithmic-cutoff} and
\eqref{eq:low-mode-variation-bound} imply
\begin{equation}\label{eq:cylindrical-transition-control}
\mathcal E_{i,n,R}^{\rm cut}
\leq
\frac{C}{(\log R)^{1/2}}\|v_n\|_{\mathcal H}^2
+o_n(1)+o_R(1),
\end{equation}
where $\mathcal E_{i,n,R}^{\rm cut}$ is the sum of the commutator
and cross terms produced by $\Delta(\eta_{i,n,R}v_n)$.  Indeed, the
high modes give the stronger factor $(\log R)^{-1}$ by
\eqref{eq:high-mode-logarithmic-cutoff}.  For the low modes,
Cauchy's inequality, the factor $(\log R)^{-1}$ carried by
$\eta'$, and \eqref{eq:low-mode-variation-bound} give the displayed
$(\log R)^{-1/2}$ bound.  This loss is immaterial and avoids any
unjustified sharp cutoff rate.
Summing the spherical harmonics in
\eqref{eq:Emden-Fowler-biharmonic} gives
\begin{equation}\label{eq:quadratic-local-energy}
\begin{aligned}
\int_\Omega|\Delta v_n|^2\,dx
\geq{}&
\sum_i\int_\Omega|\Delta v_{i,n}|^2\,dx
+(1-\delta_R)\int_\Omega|\Delta v_{0,n}|^2\,dx -\delta_R\|v_n\|_{\mathcal H}^2-o_n(1),
\end{aligned}
\end{equation}
with $\delta_R\to0$. This is the required critical logarithmic
localization.

We next localize the nonlocal form. Write
\[
d\Gamma_n(x,y)
:=\varepsilon_n^{8-\alpha}
\frac{K(x)K(y)e^{u_{\varepsilon_n}(x)+u_{\varepsilon_n}(y)}}
{|x-y|^\alpha}\,dxdy
\]
and
\[
\mathcal B_n(v,w)
:=\frac12\int_{\Omega\times\Omega}
(v(x)+v(y))(w(x)+w(y))\,d\Gamma_n(x,y).
\]
Then
$\mathcal Q_{\varepsilon_n}(v,w)
=\int_\Omega\Delta v\Delta w-\mathcal B_n(v,w)$. Let
\[
B_{i,n}=B_{R^2/\mu_{i,n}}(\xi_{i,n}),
\qquad
B_{i,n}^{\rm in}=B_{R/\mu_{i,n}}(\xi_{i,n}).
\]
The core expansion and Lemma
\ref{lem:separated-core-estimates} give
\begin{align}
\Gamma_n(B_{i,n}\times B_{j,n})
&\leq C\varepsilon_n^\alpha,
\qquad i\ne j,
\notag\\
\Gamma_n\left(
\left(\Omega\setminus\bigcup_iB_{i,n}^{\rm in}\right)
\times\Omega\right)
&\leq\delta_R+o_n(1).
\label{eq:Gamma-outer-mass}
\end{align}
The moment conditions and Lemma
\ref{lem:weighted-low-mode-control} give
\begin{equation*}
\sum_{i=1}^m
\int_{B_{i,n}}F_{\mu_{i,n},\xi_{i,n}}v_n^2\,dx
\leq C\|v_n\|_{\mathcal H}^2+o_n(1).
\end{equation*}
Cauchy's inequality with respect to $d\Gamma_n$ now shows that the
sum of the distinct-core, core--exterior, and transition errors is
bounded by
\begin{equation}\label{eq:quadratic-nonlocal-error}
|\mathcal E_{n,R}^{\rm nl}|
\leq(\delta_R+o_n(1))\|v_n\|_{\mathcal H}^2.
\end{equation}
For transition terms we also use
\eqref{eq:cylindrical-transition-control}. Combining
\eqref{eq:quadratic-local-energy} and
\eqref{eq:quadratic-nonlocal-error} proves
\eqref{eq:quadratic-localization}.

For \eqref{eq:quadratic-core-limit}, set
$x=\xi_{i,n}+z/\mu_{i,n}$. The local coefficients converge on every
fixed ball and the projected correction converges in
$C^{3,\beta}_{\rm loc}$. The part with $|z|>R$ is $o_R(1)$ by
\eqref{eq:Gamma-outer-mass}, and a distinct-core interaction is
$O(\varepsilon_n^\alpha)$. This proves
\eqref{eq:quadratic-core-limit}.
\end{proof}

\end{document}
